\chapter{Finding I --- perception is the bottleneck}
\label{ch:perception}

The single most consequential result of the reproduction is that, on the seven-step task,
the obstacle to learning is not the reinforcement-learning objective, the attention
mechanism, or the memory --- it is \emph{perception}. A front-end trained from reward alone
never learns to represent orientation, and the agent therefore cannot see the change it is
asked to report.

\section{The collapse}
The exact published network, trained end-to-end for the full $150{,}000$ episodes with the
VAE encoder learned from reward, does not reproduce the task. It converges to
\emph{always-press-at-the-deadline}: it presses at frame $5$ (the first frame a change
could appear) on every trial, scoring hit rate $1.000$ on change trials and false-alarm
rate $1.000$ on no-change trials, for an overall accuracy of $0.515$ and a reaction time
pinned at exactly five frames. It has learned the \emph{timing} of the change window but
not the \emph{discrimination} of change from no-change.

\section{The diagnosis}
Why can it not discriminate? Because to detect an orientation change one must represent
orientation, and the from-scratch front-end does not. We probed this directly: render
single Gabor patches at known orientations, encode them through the trained front-end's
$\oflat$ layer, and fit a linear decoder for the (circular) orientation. The result is
decisive (Fig.~\ref{fig:vaer2}):
\[
\Rsq_{\text{from-scratch}}=0.24,\qquad \Rsq_{\text{pretrained}}=0.97.
\]
The reward-trained front-end's features barely encode orientation; a reconstruction-
pretrained VAE's features encode it almost perfectly. Sparse reward (on the order of one
bit per episode) simply cannot teach fine orientation perception, which is exactly why the
paper pretrains the VAE as a feed-forward preprocessing step --- a step our first
reproduction had omitted.

\begin{figure}[t]
\centering
\includegraphics[width=0.52\linewidth]{fig_vae_r2.pdf}
\caption{Orientation decodability ($\Rsq$ of a linear decoder on the front-end's
$\oflat$ features) for the front-end trained from reward versus pretrained on
reconstruction. Reward alone cannot teach orientation perception; reconstruction can.}
\label{fig:vaer2}
\end{figure}

\section{The rescue}
Pretraining the VAE is cheap and decisive. We sample $\sim\!42{,}000$ grayscale
$25\times25$ patches from the environment (blanks, cues, Gabors), train the paper's exact
VAE (convolutional encoder $\to\oflat\to$ latent, transposed-convolution decoder,
reconstruction $+$ small KL) for $\sim\!12$ epochs to a reconstruction loss of $0.0025$,
and freeze the encoder into the front-end. With this single change --- \emph{nothing else
in the network or harness altered} --- the exact paper network \textbf{learns the task on
the first attempt} and climbs above the $0.5$ floor quickly, reaching hit rate $0.79$ with
a false-alarm rate of $0.08$ and a median reaction time of five frames (it presses the
instant it perceives the change).

\section{The VAE is a static, task-parameter-invariant front-end}
A subtle but useful corollary: the VAE is a \emph{per-frame, static} encoder. It
reconstructs one patch at a time and has no notion of motion or time. Task temporal
parameters --- the orientation-noise $\sigma$ and the horizon $T$ --- do not change what the
VAE sees, because the displayed orientations span the full range regardless of $\sigma$.
A single pretrained VAE therefore serves all noise levels and both horizons; there is no
need to re-pretrain per task variant. This both simplifies the pipeline and sharpens the
interpretation of the next chapter: when the long task succeeds without the pretrained VAE,
it is not because the perception problem went away, but because the long horizon admits a
different solution to it.

\section{Why this matters for the architecture choice}
This finding cleanly separates two questions that the earlier history had conflated.
``Does the architecture work?'' and ``Can the front-end see?'' are different. Many of the
collapses catalogued in Chapter~\ref{ch:zoo} on the seven-step task are, in retrospect,
perception failures wearing an architecture's clothes. The methodological upgrade is to
\emph{fix perception first} (frozen pretrained VAE) and only then compare attention/memory
wirings --- which is precisely the controlled setting of Chapters~\ref{ch:deepdive} and
\ref{ch:attention}.

A third perception result (Chapter~\ref{ch:jepa}) qualifies this: perception need not be
solved \emph{before} the architecture if the network carries a self-supervised predictive
auxiliary --- a V-JEPA temporal self-distillation loss --- trained jointly with the policy,
which learns a working front-end from raw pixels or a convolutional encoder with no
pretrained VAE.
